\emph{(i)} By Lemma \claimref{odd-extensions.volume-element-odd-dimension}, for $n$ odd $\omega$ is central and $\omega^{2}(p,q) = (-1)^{n(n-1)/2 + q}$. Then $\omega^{2}(q,p)/\omega^{2}(p,q) = (-1)^{p - q} = (-1)^{n} = -1$. The residue of $(q,p)$ is $q - p \equiv -(p - q) \pmod 8$. By Theorem \claimref{odd-extensions.doubling-or-complexification}(i) and (ii), $\omega^{2} = +1$ is the doubled type and $\omega^{2} = -1$ the complexified type, so reversal exchanges them.

\emph{(ii)} In the causal reading the time sign of $(1,3)$ is $+$ and of $(3,1)$ is $-$ (Proposition \claimref{dark-matter.towers-are-residues-single-clock-filter}: the timelike direction is the one of the minority sign). Adjoining a generator of the time sign gives $(2,3)$ and $(3,2)$; adjoining one of the other sign gives $(1,4)$ and $(4,1)$. The residues are $7, 1, 5, 3$. Lemma \claimref{odd-extensions.volume-element-odd-dimension} at $n = 5$ gives $\omega^{2} = (-1)^{q}$: $+1$ for $(1,4)$ and $(3,2)$, $-1$ for $(2,3)$ and $(4,1)$, and Theorem \claimref{odd-extensions.doubling-or-complexification} gives the types. Reversing $(1,4)$ gives $(4,1)$ and reversing $(2,3)$ gives $(3,2)$, so mirror $+ X$ is the reversal of this world $+ X$ for $X$ = space and time; $5 + 3 \equiv 7 + 1 \equiv 0 \pmod 8$. Both shadows have $q$ odd; adjoining a generator of square $-1$ makes $q$ even and $\omega^{2} = +1$, adjoining one of square $+1$ leaves $q$ odd and $\omega^{2} = -1$, in either tower. This is the computation in the proof of Corollary \claimref{odd-extensions.four-extensions-oriented-pair}, whose algebras and residues are unchanged; only the labels of the mirror's two rows are exchanged.

\emph{(iii)} In the uniform reading the mirror's space generator has square $-1$, giving $(3,2)$ at residue $1$, and its time generator square $+1$, giving $(4,1)$ at residue $3$, as Corollary \claimref{odd-extensions.four-extensions-oriented-pair} lists. Then $5 + 1 \equiv 6$ and $7 + 3 \equiv 2 \pmod 8$, neither $0$, and the mirror's extension of each kind is not the reversal of ours. With time the plus sign, $(3,1)$ has three time directions. It still has $\min(p,q) = 1$, so it is causal by the Proposition's definition; what the uniform reading contradicts is that Proposition's reading of the single minority direction as the timelike one.

\emph{(iv)} Theorem \claimref{rule2.where-obstruction-sits-commuting-square}(a): for $n$ even and $p - q \equiv 2 \pmod 4$ the obstructed extension is $A_{s} = \Cl(p, q{+}1)$, the one by a generator of square $-1$, and the extension by a generator of square $+1$ has complex type. Both $(1,3)$ and $(3,1)$ have $p - q \equiv 2 \pmod 4$. For $(1,3)$ the square $-1$ generator is spacelike in either reading, giving $\Cl(1,4)$ at residue $5$, which is Corollary \claimref{rule2.which-residues-obstruct-which-way}. For $(3,1)$ in the causal reading the square $-1$ generator is timelike: the obstructed extension $\Cl(3,2)$ is a second time direction, and the spacelike extension is $\Cl(4,1)$, of complex type, simple, residue $3$. The corollary's sentence on residue $2$ is therefore a statement about the extension by a generator of square $-1$, the theorem's $A_{s}$; it names the mirror's Kaluza-Klein direction only in the uniform reading.

\emph{(v)} By Proposition \claimref{odd-extensions.two-towers-one-curved-universe}(i), for $\Lambda > 0$ our $I^{\perp}$ is $(1,4)$ and the mirror's, its reversal, is $(4,1)$; the added unit vector $e_{V}$ has $Q(e_{V}) = -1$ in ours and $Q(e_{V}) = +1$ in the mirror's. In the causal reading these are the space signs of the two towers, so both add a spatial direction; the types are those of (ii). In the uniform reading the mirror's $Q(e_{V}) = +1$ is the time sign.

\emph{The convention of T14 and of the Rule 2 plane.}  Proposition \claimref{dark-matter.nesting}(i) builds the residue-five tower as this world plus a fifth spatial generator and states that the mirror chain is the same with the signatures reversed, which is the causal reading of (ii). Proposition \claimref{rule2.rule-2-plane}(ii) reads $\Cl(3,1)$ as $\Cl(2,1)$ plus a third spatial direction with $\nu^{2} = +1$, the plane's time $e_{t}$ having square $-1$: again the minority sign carries time. Theorem \claimref{rule2.where-obstruction-sits-commuting-square} instead names $A_{s} = \Cl(p,q{+}1)$ spacelike for every $(p,q)$, which for the mirror is the uniform reading, so the Rule 2 note uses both. Proposition \claimref{dark-matter.towers-are-residues-single-clock-filter}'s phrase "time $=$ plus" labels this world $(1,3)$, where plus is the minority sign, and does not conflict with the causal reading.

The check recomputes $\omega^{2}$ and the centrality of $\omega$ for every signature of odd dimension up to nine, the central idempotents of the doubled cases, the table of (ii) and its failure under the uniform reading, the obstructed extension for every even signature up to dimension eight, and the de Sitter vector on rational instances of the frozen point in the conformal-model form $I = n_{\infty} + (\Lambda/6)\,n_{0}$, in the tractor signatures $(2,4)$ and $(4,2)$, reading each world's space sign as the majority sign of its Lorentz frame rather than from the convention being tested.
